\newpage
\section{Glosario}
\subsection*{Capítulo 1. Planteamiento del proyecto}
\begin{description}
    \item[Plaga forestal:] Organismo o conjunto de organismos que afectan negativamente la salud, crecimiento o supervivencia del arbolado.
    \item[Insectos descortezadores:] Escarabajos que viven y se desarrollan debajo de la corteza de los árboles, alimentándose de tejidos conductores y provocando debilitamiento o muerte del árbol.
    \item[Detección temprana:] Identificación de señales iniciales de una afectación antes de que el daño sea severo o irreversible.
    \item[Arbolado:] Conjunto de árboles presentes en una zona determinada, como un bosque, parque o área verde urbana.
    \item[Monitoreo:] Proceso de observación, registro y seguimiento continuo del estado de un fenómeno, sistema o recurso natural.
    \item[Visión por computadora:] Área de la inteligencia artificial que permite a los sistemas analizar e interpretar imágenes o videos para extraer información útil.
    \item[Inteligencia artificial:] Campo de la computación orientado al desarrollo de sistemas capaces de realizar tareas que normalmente requieren inteligencia humana.
    \item[Procesamiento de imágenes:] Conjunto de técnicas utilizadas para modificar, mejorar o analizar imágenes digitales.
    \item[Aplicación móvil:] Software diseñado para ejecutarse en dispositivos móviles, como teléfonos inteligentes o tabletas.
    \item[Dataset:] Conjunto de datos organizado que se utiliza para entrenar, validar o probar un modelo computacional.
    \item[Scrum:] Metodología ágil de gestión de proyectos basada en ciclos iterativos llamados sprints.
    \item[Sprint:] Periodo corto de trabajo dentro de Scrum en el que se desarrollan tareas específicas del proyecto.
    \item[Product Backlog:] Lista priorizada de requerimientos, funcionalidades o tareas pendientes de un proyecto ágil.
    \item[Sprint Backlog:] Conjunto de tareas seleccionadas para desarrollarse durante un sprint específico.
\end{description}
\subsection*{Capítulo 2. Contexto y caracterización del caso de estudio}
\begin{description}
    \item[Bosque urbano:] Área verde con presencia significativa de árboles ubicada dentro o cerca de una zona urbana.
    \item[Servicios ecosistémicos:] Beneficios que los ecosistemas proporcionan a la sociedad, como regulación de temperatura, mejora de calidad del aire y espacios de recreación.
    \item[Infestación:] Presencia y desarrollo de organismos dañinos, como insectos descortezadores, dentro o sobre un árbol.
    \item[Hospedero:] Árbol o planta que sirve como fuente de alimento, refugio o sitio de reproducción para un organismo.
    \item[Agente biótico:] Organismo vivo que influye sobre otro organismo o sobre un ecosistema.
    \item[Coleóptero:] Insecto perteneciente al orden Coleoptera, caracterizado por tener alas anteriores endurecidas llamadas élitros.
    \item[Curculionidae:] Familia de insectos coleópteros a la que pertenecen diversos escarabajos, incluidos algunos descortezadores.
    \item[Dendroctonus:] Género de insectos descortezadores de alta relevancia forestal debido a su capacidad para causar la muerte de árboles.
    \item[Ips:] Género de insectos descortezadores asociado al ataque de árboles debilitados o estresados.
    \item[Hylesinus sp.:] Género de descortezador asociado a afectaciones en fresnos, especialmente en ramas y tallos.
    \item[Phloeosinus sp.:] Género de descortezador asociado a afectaciones en cedros y otras especies de la familia Cupressaceae.
    \item[Fraxinus uhdei:] Nombre científico del fresno, especie arbórea presente en el Bosque de Chapultepec.
    \item[Cupressus lusitanica:] Nombre científico del cedro blanco, especie susceptible a infestaciones por descortezadores.
    \item[Corteza:] Capa externa del tronco y ramas del árbol que protege los tejidos internos.
    \item[Tejido subcortical:] Región ubicada debajo de la corteza del árbol, donde pueden desarrollarse galerías producidas por descortezadores.
    \item[Floema:] Tejido vegetal encargado de transportar nutrientes producidos por la fotosíntesis hacia distintas partes del árbol.
    \item[Xilema:] Tejido vegetal encargado de transportar agua y minerales desde las raíces hacia el resto del árbol.
    \item[Resina:] Sustancia viscosa producida por algunos árboles como mecanismo de defensa ante heridas o ataques de insectos.
    \item[Grumos de resina:] Acumulaciones visibles de resina sobre la corteza que pueden indicar perforaciones o ataques de insectos.
    \item[Galerías subcorticales:] Túneles excavados por insectos descortezadores debajo de la corteza durante su alimentación o reproducción.
    \item[Orificios de entrada:] Perforaciones realizadas por los insectos para ingresar al interior de la corteza.
    \item[Orificios de salida:] Perforaciones por donde emergen los insectos adultos después de completar su desarrollo.
    \item[Estrés hídrico:] Condición en la que un árbol no dispone de suficiente agua para mantener sus funciones fisiológicas normales.
    \item[Colonización:] Proceso mediante el cual los insectos descortezadores ingresan y se establecen en un árbol hospedero.
    \item[Follaje:] Conjunto de hojas o acículas de un árbol, cuya coloración puede indicar su estado sanitario.
    \item[Dosel:] Capa superior formada por las copas de los árboles en un bosque.
    \item[Infestación activa:] Estado en el que los insectos se encuentran vivos, alimentándose o reproduciéndose dentro del árbol.
    \item[Fase de infestación:] Etapa del proceso de daño del árbol, distinguible por señales visuales como cambios de color en follaje o corteza.
\end{description}
\subsection*{Capítulo 3. Delimitación computacional del problema}
\begin{description}
    \item[Información de entrada:] Datos que recibe un sistema para realizar un procesamiento; en este proyecto corresponde principalmente a imágenes de árboles.
    \item[Imagen digital:] Representación visual compuesta por píxeles que puede ser procesada por una computadora.
    \item[Píxel:] Unidad mínima de una imagen digital, asociada a un valor de color o intensidad.
    \item[Clasificación de imágenes:] Tarea de visión por computadora que asigna una imagen a una categoría o clase determinada.
    \item[Detección de objetos:] Tarea de visión por computadora que localiza e identifica objetos específicos dentro de una imagen.
    \item[Segmentación de imágenes:] Técnica que divide una imagen en regiones o zonas significativas para su análisis.
    \item[Características visuales:] Rasgos presentes en una imagen, como color, textura, forma o patrones, que pueden ser utilizados para su análisis.
    \item[Extracción de características:] Proceso mediante el cual se identifican y obtienen rasgos relevantes de una imagen.
    \item[Modelo computacional:] Representación matemática o algorítmica capaz de procesar datos y producir resultados.
    \item[Flujo de trabajo:] Secuencia organizada de pasos necesarios para desarrollar, entrenar, evaluar o implementar un sistema.
\end{description}
\subsection*{Capítulo 4. Fundamentación del enfoque de aprendizaje profundo}
\begin{description}
    \item[Aprendizaje automático:] Rama de la inteligencia artificial que permite a los sistemas aprender patrones a partir de datos.
    \item[Aprendizaje supervisado:] Tipo de aprendizaje automático en el que el modelo se entrena con datos previamente etiquetados.
    \item[Clasificación:] Tarea de aprendizaje automático que consiste en asignar una entrada a una de varias clases definidas.
    \item[Clase:] Categoría o etiqueta que representa el resultado esperado de una predicción.
    \item[Aprendizaje profundo:] Subcampo del aprendizaje automático basado en redes neuronales con múltiples capas.
    \item[Red neuronal:] Modelo computacional inspirado en el funcionamiento del cerebro, formado por unidades llamadas neuronas artificiales.
    \item[Red neuronal convolucional:] Tipo de red neuronal especializada en el análisis de imágenes mediante operaciones de convolución.
    \item[Datos anotados:] Datos que incluyen etiquetas o información adicional que indica la respuesta esperada para el entrenamiento de un modelo.
    \item[Entrenamiento:] Proceso mediante el cual un modelo ajusta sus parámetros para aprender patrones a partir de datos.
    \item[Validación:] Evaluación del modelo durante el desarrollo para medir su desempeño y ajustar decisiones de entrenamiento.
    \item[Prueba:] Evaluación final del modelo con datos que no fueron utilizados durante el entrenamiento.
    \item[Generalización:] Capacidad de un modelo para producir resultados correctos ante datos nuevos o no vistos previamente.
    \item[Sobreajuste:] Fenómeno en el que un modelo aprende demasiado los datos de entrenamiento y pierde capacidad de generalizar.
    \item[Precisión:] Métrica que indica la proporción de predicciones correctas realizadas por un modelo.
    \item[Predicción:] Resultado generado por un modelo al analizar una entrada.
    \item[Inferencia:] Proceso de utilizar un modelo entrenado para generar predicciones sobre nuevos datos.
\end{description}
\subsection*{Capítulo 5. Análisis comparativo de arquitecturas y selección del modelo}
\begin{description}
    \item[Arquitectura de modelo:] Estructura interna de un modelo de aprendizaje profundo, definida por sus capas, conexiones y operaciones.
    \item[ResNet-50:] Arquitectura de red neuronal profunda basada en conexiones residuales, utilizada en clasificación de imágenes.
    \item[MobileNetV3:] Arquitectura de red neuronal optimizada para dispositivos móviles y sistemas con recursos limitados.
    \item[EfficientNet:] Familia de modelos de aprendizaje profundo que busca equilibrar precisión y eficiencia computacional.
    \item[EfficientNet-B3:] Variante de EfficientNet con mayor capacidad que modelos más pequeños de la misma familia, utilizada para clasificación de imágenes.
    \item[MobileViT:] Arquitectura que combina redes convolucionales con mecanismos tipo Transformer para tareas visuales.
    \item[Transfer learning:] Técnica que reutiliza un modelo previamente entrenado para adaptarlo a una nueva tarea.
    \item[Fine-tuning:] Ajuste fino de un modelo preentrenado mediante entrenamiento adicional con datos específicos del problema.
    \item[Backbone:] Red base de una arquitectura que extrae características generales de las imágenes.
    \item[Modelo preentrenado:] Modelo que ya fue entrenado previamente con un conjunto de datos amplio y puede reutilizarse en otra tarea.
    \item[ImageNet:] Conjunto de datos ampliamente utilizado para entrenar y evaluar modelos de clasificación de imágenes.
    \item[Convolución:] Operación matemática que permite detectar patrones locales en una imagen, como bordes, texturas o formas.
    \item[Batch Normalization:] Técnica que normaliza valores intermedios dentro de una red neuronal para estabilizar y acelerar el entrenamiento.
    \item[ReLU:] Función de activación que transforma valores negativos en cero y conserva valores positivos.
    \item[Pooling:] Operación que reduce el tamaño de los mapas de características conservando información relevante.
    \item[Flatten:] Operación que convierte una matriz o mapa de características en un vector.
    \item[Capa densa:] Capa neuronal en la que cada neurona está conectada con todas las neuronas de la capa anterior.
    \item[Softmax:] Función que convierte los valores de salida de un modelo en probabilidades asociadas a cada clase.
    \item[Parámetros del modelo:] Valores internos que el modelo ajusta durante el entrenamiento para aprender patrones.
    \item[Escalado compuesto:] Estrategia utilizada en EfficientNet para aumentar de forma balanceada la profundidad, ancho y resolución del modelo.
\end{description}
\subsection*{Capítulo 6. Conjunto de datos y preprocesamiento}
\begin{description}
    \item[Recopilación de datos:] Proceso de obtener imágenes o información necesaria para construir el conjunto de datos del proyecto.
    \item[Conjunto de datos:] Colección organizada de ejemplos utilizados para entrenar, validar y probar un modelo.
    \item[Etiquetado de imágenes:] Proceso de asignar una clase o categoría a cada imagen del conjunto de datos.
    \item[Preprocesamiento:] Conjunto de operaciones aplicadas a los datos antes de utilizarlos en el entrenamiento de un modelo.
    \item[Normalización:] Ajuste de los valores de una imagen para que se encuentren dentro de un rango común y faciliten el entrenamiento.
    \item[Redimensionamiento:] Modificación del tamaño de una imagen para adaptarla a los requerimientos de entrada del modelo.
    \item[Aumentación de datos:] Técnica que genera variaciones de imágenes existentes para incrementar la diversidad del conjunto de datos.
    \item[Conjunto de entrenamiento:] Parte del dataset utilizada para que el modelo aprenda patrones.
    \item[Conjunto de validación:] Parte del dataset utilizada para evaluar el comportamiento del modelo durante el desarrollo.
    \item[Conjunto de prueba:] Parte del dataset utilizada para medir el desempeño final del modelo.
    \item[Análisis cromático:] Evaluación de los colores presentes en una imagen para identificar patrones asociados al estado del árbol.
    \item[Coloración del follaje:] Cambio visual en las hojas o copa del árbol que puede indicar una fase de deterioro o infestación.
\end{description}
\subsection*{Capítulo 7. Análisis del sistema}
\begin{description}
    \item[Requerimiento funcional:] Función o comportamiento específico que el sistema debe realizar.
    \item[Requerimiento no funcional:] Condición de calidad que debe cumplir el sistema, como rendimiento, seguridad, disponibilidad o usabilidad.
    \item[Regla de negocio:] Condición o restricción que define cómo debe comportarse el sistema según el contexto del proyecto.
    \item[Actor del sistema:] Usuario, entidad o componente externo que interactúa con el sistema.
    \item[Métrica de evaluación:] Criterio cuantitativo utilizado para medir el desempeño del modelo o del sistema.
    \item[Criterio de aceptación:] Condición que debe cumplirse para considerar válida una funcionalidad o entregable.
    \item[Riesgo técnico:] Posible problema relacionado con la implementación, funcionamiento o integración del sistema.
    \item[Mitigación:] Acción planificada para reducir la probabilidad o impacto de un riesgo.
    \item[Prueba funcional:] Evaluación que verifica si una función del sistema opera correctamente.
    \item[Prueba de integración:] Evaluación que verifica la comunicación y funcionamiento conjunto de varios módulos del sistema.
    \item[Infraestructura en la nube:] Conjunto de recursos computacionales ofrecidos por un proveedor de nube para almacenamiento, procesamiento o despliegue.
    \item[AWS:] Plataforma de servicios en la nube de Amazon utilizada para alojar, procesar o administrar componentes del sistema.
    \item[Azure:] Plataforma de servicios en la nube de Microsoft considerada como alternativa de infraestructura.
    \item[GCP:] Plataforma de servicios en la nube de Google considerada como alternativa de infraestructura.
    \item[Amazon SageMaker:] Servicio de AWS orientado al desarrollo, entrenamiento y despliegue de modelos de aprendizaje automático.
    \item[AWS Lambda:] Servicio serverless de AWS que permite ejecutar código sin administrar servidores directamente.
    \item[Amazon RDS:] Servicio de AWS para administrar bases de datos relacionales.
    \item[PostgreSQL:] Sistema de gestión de bases de datos relacional de código abierto.
    \item[Serverless:] Modelo de computación en la nube en el que el proveedor administra la infraestructura y el usuario solo define la lógica de ejecución.
    \item[Contenedor:] Unidad de software que empaqueta una aplicación y sus dependencias para ejecutarse de forma portable.
    \item[Capa gratuita:] Conjunto de servicios o recursos ofrecidos sin costo limitado por un proveedor de nube.
    \item[AWS Pricing Calculator:] Herramienta de AWS utilizada para estimar costos de infraestructura y servicios en la nube.
\end{description}
\subsection*{Capítulo 8. Diseño del sistema}
\begin{description}
    \item[Arquitectura del sistema:] Organización general de los componentes que conforman el sistema y sus relaciones.
    \item[Backend:] Parte del sistema encargada de la lógica de negocio, procesamiento, comunicación con bases de datos y servicios.
    \item[Frontend:] Parte visual del sistema con la que interactúa el usuario.
    \item[API:] Interfaz que permite la comunicación entre diferentes componentes de software.
    \item[Clean Architecture:] Patrón de diseño que organiza el software en capas independientes para mejorar mantenibilidad y separación de responsabilidades.
    \item[CI/CD:] Prácticas de integración continua y despliegue continuo utilizadas para automatizar pruebas, construcción y liberación de software.
    \item[Integración continua:] Práctica que consiste en incorporar frecuentemente cambios de código y verificarlos mediante procesos automáticos.
    \item[Despliegue continuo:] Práctica que automatiza la publicación de nuevas versiones del sistema.
    \item[Observabilidad:] Capacidad de conocer el estado interno de un sistema mediante métricas, registros y trazas.
    \item[Monitoreo centralizado:] Supervisión unificada de eventos, errores y métricas de los componentes del sistema.
    \item[Offline-First:] Estrategia de diseño que permite que una aplicación funcione aun cuando no exista conexión a internet.
    \item[Sincronización:] Proceso mediante el cual los datos locales y remotos se actualizan para mantener consistencia.
    \item[Inferencia local:] Ejecución del modelo de clasificación directamente en el dispositivo móvil.
    \item[SPA:] Aplicación web de una sola página que actualiza su contenido dinámicamente sin recargar completamente el sitio.
    \item[Autenticación:] Proceso mediante el cual se verifica la identidad de un usuario.
    \item[Control de acceso:] Mecanismo que determina qué acciones o recursos puede utilizar un usuario dentro del sistema.
    \item[Dashboard:] Interfaz que presenta información resumida, indicadores o accesos principales del sistema.
    \item[UI:] Interfaz de usuario, es decir, los elementos visuales con los que interactúa una persona.
    \item[UX:] Experiencia de usuario, relacionada con la facilidad, claridad y satisfacción al utilizar el sistema.
    \item[Caso de uso:] Descripción de una interacción entre un actor y el sistema para lograr un objetivo específico.
    \item[Modelo entidad-relación:] Representación conceptual de las entidades, atributos y relaciones de una base de datos.
    \item[Esquema relacional:] Organización lógica de una base de datos mediante tablas, columnas y relaciones.
    \item[Diccionario de datos:] Documento que describe las tablas, campos, tipos de datos y restricciones de una base de datos.
    \item[Entidad:] Objeto o concepto del mundo real representado dentro de una base de datos.
    \item[Atributo:] Propiedad o característica de una entidad.
    \item[Llave primaria:] Atributo que identifica de forma única cada registro de una tabla.
    \item[Llave foránea:] Atributo que establece una relación entre dos tablas.
    \item[Persistencia:] Capacidad de almacenar datos de forma duradera para consultarlos posteriormente.
    \item[Persistencia políglota:] Uso de diferentes tecnologías de almacenamiento según el tipo o necesidad de los datos.
    \item[Análisis cromático:] Proceso complementario que evalúa la información de color de las imágenes para apoyar la clasificación del estado del árbol.
    \item[Registro:] Entrada almacenada en el sistema que contiene información asociada a un árbol, análisis o usuario.
\end{description}
\subsection*{Capítulo 9. Pruebas y validación del sistema}
\begin{description}
    \item[Benchmark:] Proceso sistemático de evaluación y comparación cuantitativa del rendimiento de un sistema frente a umbrales predefinidos.
    \item[PSS (Proportional Set Size):] Medida de memoria física real atribuible a un proceso, calculada sumando la memoria privada más la fracción proporcional de memoria compartida con otros procesos.
    \item[Java Heap:] Región de memoria administrada automáticamente por la máquina virtual de Android (ART) para instancias de objetos de la aplicación.
    \item[Native Heap:] Espacio de memoria asignado dinámicamente fuera de la máquina virtual gestionada, comúnmente utilizado por bibliotecas en C y C++ como TensorFlow Lite o MapLibre GL Native.
    \item[Low Memory Killer (LMK):] Mecanismo del kernel de Linux en Android encargado de terminar selectivamente procesos cuando la memoria física disponible alcanza niveles críticos.
    \item[Janky Frames:] Cuadros de la interfaz de usuario cuyo tiempo de renderizado supera el presupuesto permitido por la tasa de refresco de la pantalla (e.g., más de 16.6~ms a 60~Hz o más de 8.3~ms a 120~Hz), produciendo tirones visibles.
    \item[Android Debug Bridge (ADB):] Herramienta de línea de comandos que permite la comunicación, depuración y control de un dispositivo Android conectado.
    \item[AVD (Android Virtual Device):] Configuración que emula las características de hardware y software de un dispositivo Android físico dentro del entorno de desarrollo.
    \item[Firebase Test Lab:] Infraestructura basada en la nube de Google que permite ejecutar pruebas automatizadas en matrices de dispositivos móviles físicos y virtuales.
    \item[Fixed-Focus:] Óptica de cámara con distancia de enfoque fija e invariable, usualmente calibrada a la distancia hiperfocal sin posibilidad de enfoque macro cercano.
    \item[Write-Ahead Logging (WAL):] Mecanismo de persistencia transaccional en bases de datos relacionales donde las modificaciones se registran secuencialmente en un archivo de bitácora antes de consolidarse en el archivo principal.
    \item[OpenGL ES:] Especificación estándar de interfaz de programación de aplicaciones para el renderizado de gráficos vectoriales 2D y 3D en dispositivos embebidos y móviles.
    \item[Fine-Tuning:] Técnica de ajuste fino en aprendizaje por transferencia donde se descongelan y reentrenan capas convolucionales superiores a una tasa de aprendizaje reducida.
    \item[Early Stopping:] Criterio de parada temprana que interrumpe el entrenamiento de una red neuronal cuando el desempeño en el conjunto de validación deja de mejorar, evitando el sobreajuste.
    \item[Data Leakage:] Fuga de información que ocurre cuando datos del conjunto de prueba o validación influyen inadvertidamente en el entrenamiento del modelo.
    \item[F-beta Score:] Métrica de evaluación que pondera la importancia relativa entre la precisión y la exhaustividad mediante un factor de ponderación $\beta$.
    \item[Curva Precision-Recall:] Gráfica que ilustra el balance entre la precisión y la exhaustividad de un clasificador binario a través de diversos umbrales de decisión.
    \item[Paridad numérica:] Grado de consistencia matemática en las predicciones generadas por un modelo al ser ejecutado en diferentes entornos o motores de inferencia (e.g., Keras frente a TensorFlow Lite).
\end{description}







